\chapter{The Web via Emscripten}
\label{ch:web-emscripten}

CNA's Web target compiles the same C\texttt{++}23 codebase through Emscripten rather than a
native toolchain, and this chapter covers the specific build-system machinery that makes that
possible, along with an honest account of exactly how far it has actually been verified.

\section{Build-system adjustments specific to Emscripten}

Two adjustments apply project-wide before any backend-specific configuration happens.
Emscripten disables C\texttt{++} exception unwinding by default, for code-size reasons; CNA
force-re-enables it globally, since exceptions are used throughout the codebase (every
\texttt{NotImplementedException}, every \texttt{ArgumentOutOfRangeException} this book has
cited depends on real unwinding working). And \texttt{EASYGL} (targeting WebGL2, which is
GLES~3.0) is the default backend under Emscripten, the same default it holds on Linux ---
while \texttt{CANVAS} (Chapter~\ref{ch:canvas-ascii-backends}) is hard-gated the other
direction, refusing to configure at all outside Emscripten, since HTML \texttt{<canvas>} has
no meaning anywhere else.

Several test categories are excluded specifically under Emscripten, each for a concrete
reason rather than a blanket ``skip on web'' rule: pixel/integration test suites for every
backend's own CMake test registration are excluded here, and a handful of process-spawning
tests (a two-process networking loopback test, a gamer-services-dispatcher hang-regression
test, an audio mixer test, a glTF-tool test) are excluded specifically because there is no
real multi-process spawning available inside a single Node.js/WebAssembly module.

\section{Making the shared test binary actually exit, and actually yield}

Two further Emscripten-specific link flags apply only to the shared \texttt{CnaTests} binary,
each solving a distinct, concrete problem. Emscripten's default behavior keeps its JavaScript
runtime alive after \texttt{main()} returns, so \texttt{node CnaTests.js} would otherwise never
exit on its own --- a runtime-exit flag corrects this specifically for the test binary.
Separately, the system-link networking loopback tests need to actually yield control back to
Node's own event loop mid-test (so a WebSocket handshake, for instance, has a chance to
complete) --- ordinary synchronous, single-threaded Emscripten execution cannot yield at all
without help, confirmed directly: a real one-second sleep loop never let a WebSocket handshake
finish without Emscripten's Asyncify transformation enabled, so Asyncify is applied, scoped
specifically to this one test binary rather than the whole project.

\section{What has actually been built and run, versus what has not}

This is the section most worth reading carefully, because the picture differs sharply by
backend. The \texttt{CANVAS} backend has a real, confirmed success story: a genuine
\texttt{emcmake}/\texttt{emcc} (version 6.0.2) configure-and-build succeeded, producing a real
build directory for the first time in this project's history, and \texttt{CnaTests.js}
genuinely links and runs under \texttt{node}, with a real, backend-agnostic GoogleTest suite
passing in full.

The default \texttt{EASYGL}/WebGL2 path tells a different, more honest story: at the time of
writing, no Emscripten build directory has ever existed for \texttt{EASYGL} at all, and no
automated workflow builds it under Emscripten either. Every specific claim in this project's
own Web/Emscripten graphics-limitations documentation about \texttt{EASYGL} behavior under
WebGL2 --- geometry and tessellation shader absence, anisotropic-filtering extension gating,
whether a WebGL1 fallback exists --- is explicitly labeled a \emph{design-time expectation},
not something empirically verified, precisely because the build that would let someone verify
it has never actually been produced.

\section{Why even a successful build cannot be pixel-verified here}

Even where a build genuinely succeeds and a test binary genuinely runs, there is a hard
ceiling on what it can prove: \texttt{node} provides no WebGL context and no
\texttt{CanvasRenderingContext2D}/DOM at all --- confirmed directly, since
\texttt{SDL\_Init(SDL\_INIT\_VIDEO)} itself throws under Emscripten running inside
\texttt{node}. This is the same ceiling Chapter~\ref{ch:canvas-ascii-backends} already
described for \texttt{CANVAS} specifically: any test suite that constructs a real
\cnaclass{GraphicsDevice} simply cannot execute in this dev loop at all, real browser or not,
which is exactly why \texttt{docs/canvas-backend.md}'s own manual browser verification
checklist --- covering rotation and origin math, blend modes, and wrap/mirror sampling among
thirteen total items --- is left entirely unchecked, deferred to a human running the game in
an actual browser via a tool like \texttt{emrun}, rather than claimed as done.

One further, genuinely interesting piece of real code is worth knowing about even though it
has never been exercised against a real browser: \texttt{EASYGL}'s own backend implementation
contains real, reviewed WebGL-context-loss handling --- registering for the browser's own
\texttt{WEBGL\_lose\_context} extension and its \texttt{webglcontextlost}/%
\texttt{webglcontextrestored} events. This code path exists and has been read and reviewed,
but --- consistent with everything else in this section --- has never actually run against a
real browser; only the ordinary desktop code path (a synchronous SDL GL context
destroy/recreate) has ever actually executed.

\section{The honest summary}

Web support in CNA is real in the sense that a specific, narrow slice of it (\texttt{CANVAS},
building and linking under Node) has been genuinely produced and run; it is aspirational in
the sense that the backend most users would actually reach for (\texttt{EASYGL}/WebGL2) has
never been built under Emscripten at all in this project's own history, and no claim about
its WebGL2-specific behavior should be read as more than a design-time expectation until
someone actually produces that build and checks it in a real browser.
